%%=============================================================================
%% Methodologie
%%=============================================================================

\chapter{\IfLanguageName{dutch}{Methodologie}{Methodology}}\label{ch:methodologie}

Voor het bekomen van een gewenst en bruikbaar eindresultaat werd een stapsgewijze aanpak gehanteerd. In deze methode worden vier afzonderlijke fasen doorlopen, waarbij elke fase voortbouwt op de resultaten van de vorige fase. Doordat het onderzoek past in een overkoepelend onderzoekstraject naar een duurzame energiemix voor de ET zijn deze stappen niet stapsgewijs en worden deze deels over elkaar gedaan. Zo kunnen de vereisten voor een goede analyse veranderen nadat de initiele vereistenanalyse is afgerond. Hierdoor kan de tool die wordt ontwikkeld gaandeweg worden bijgestuurd zodat het eindresultaat zo optimaal mogelijk aansluit bij de behoeften van de stakeholders. Dit principe van iteratie en continue feedback is essentieel voor het leveren van een kwalitatief en bruikbaar eindresultaat.

\section{Stap 1: Vereistenanalyse}
Deze fase beoogt het identificeren van de vereisten en beperkingen specifiek voor de ET locatie in de EMR. Dit door eerst interviews af te nemen met stakeholders (ET-onderzoekers) om de benodigde datasets (bv. Natura 2000-gebieden, landbouwzones, stortplaatsen) en criteria te bepalen per geplande energiebron. Hierna volgt een analyse van de samenstelling en bruikbaarheid van bestaande aangeleverde datasets. Met de verkregen informatie kunnen ontbrekende datasets geïdentificeerd worden. Op basis van de informatie kan ook bepaald worden wat de tool moet kunnen.

\section{Stap 2: Vergelijking van bestaande tools}
Na het bepalen van de vereisten van de onderzoekers kan de meest geschikte software of bibliotheek als basis voor verdere ontwikkeling worden bepaald. Eerst worden bestaande software en libraries met hun functionaliteit, gebruiksgemak, toegankelijkheid en documentatie geidentificeerd in een literatuurstudie. Op basis hiervan kunnen de meest geschikte tools worden bepaald en samengebracht in een short-list. Deze lijst kan dan verder verkleind worden op basis van secundaire criteria totdat het meest geschikte startpunt is bepaald.

\section{Stap 3: Ontwikkeling van de Tool (Proof of Concept (PoC))}
In de voorlaatste fase wordt de geselecteerde tool uitgebreid om GIS-datasets te analyseren volgens de wensen van de onderzoekers. Er zal een aanpak bepaald worden om de analyze correct en efficiënt te laten verlopen. De bekomen analyse zal dan volgens de wensen van de gebruiker gerapporteerd worden. Hierbij zal voorkeur worden gegeven aan het werken met open standaarden. Er zal ook rekening gehouden worden met de best-practices van de software ontwikkeling, zo zal er zo veel mogelijk generiek en uitbreidbaar gewerkt worden. De code zal zo veel mogelijk opgedeeld worden om een snelle aanpassing van de aanpak mogelijk te maken.

\section{Stap 4: Analyseren van de risico's}
Finaal worden potentiële risico's geïdentificeerd en ingeschat voor de bekomen en toekomstige resultaten. Hierin wordt onderzocht of er blind-spots of valkuilen zitten in de bekomen resultaten. Deze mogelijke fouten kunnen het gevolg zijn van ontbrekende data, technische beperkingen of onvolledige eisen. Met deze laatste stap worden de bekomen resultaten zo transparant mogelijk gecommuniceerd naar de stakeholders.

%% TODO: In dit hoofstuk geef je een korte toelichting over hoe je te werk bent
%% gegaan. Verdeel je onderzoek in grote fasen, en licht in elke fase toe wat
%% de doelstelling was, welke deliverables daar uit gekomen zijn, en welke
%% onderzoeksmethoden je daarbij toegepast hebt. Verantwoord waarom je
%% op deze manier te werk gegaan bent.
%% 
%% Voorbeelden van zulke fasen zijn: literatuurstudie, opstellen van een
%% requirements-analyse, opstellen long-list (bij vergelijkende studie),
%% selectie van geschikte tools (bij vergelijkende studie, "short-list"),
%% opzetten testopstelling/PoC, uitvoeren testen en verzamelen
%% van resultaten, analyse van resultaten, ...
%%
%% !!!!! LET OP !!!!!
%%
%% Het is uitdrukkelijk NIET de bedoeling dat je het grootste deel van de corpus
%% van je bachelorproef in dit hoofstuk verwerkt! Dit hoofdstuk is eerder een
%% kort overzicht van je plan van aanpak.
%%
%% Maak voor elke fase (behalve het literatuuronderzoek) een NIEUW HOOFDSTUK aan
%% en geef het een gepaste titel.


